\chapter{The Immune System}

\section{Complement system}

\textit{This section is based primarily on \cite{DaveComplementSystem} and the week 105 lectures from Kenny McKeegan.}

The complement system is a part of the innate immune system based on proteins. The proteins that play a role in the complement system are named C1, C2, \ldots C9 after the order in which they were discovered. Each of the complement system proteins comprises two halves: an \href{https://en.wikipedia.org/wiki/Anaphylatoxin}{\textbf{anaphylatoxin}} part (labelled with an `a'), which recruits neutrophils and other phagocytes to the site of infection, and a binding part, labelled with a `b' that binds to the surface of a pathogen. 

The general idea of the complement system is that one of the binding proteins (e.g. C1b) will bind to the surface of a pathogen, which will cause the binding and cleavage of further proteins, causing the binding and cleavage of further proteins, resulting in a \href{https://en.wikipedia.org/wiki/Complement_system}{\textbf{complement cascade}}. There are three named kinds of complement cascades covered in the following subsections, but they have many similarities. In particular, all three cascades --- at some point in their execution --- produce a C3b (often via a \href{https://en.wikipedia.org/wiki/C3-convertase}{\textbf{C3 convertase}}: a protein that cleaves C3 into C3a and C3b). Once C3b is produced, it binds to the pathogen and \textbf{opsonizes} it, among other things, it and the cascades become the same.

In particular, once C3b has bonded to the pathogen, the cascade continues as follows:

\begin{itemize}
    \item \textbf{Factor B} in the plasma binds to C3b, which is converted by \textbf{Factor D} into C3bBb on the pathogen surface.
    \item This C3bBb complex is stabilized by \href{https://en.wikipedia.org/wiki/Properdin}{\textbf{properdin}} and acts as a C3 convertase causing a positive feedback loop.
    \item Some of the C3bBb is bonded by another C3b to form C3b\({}_2\)Bb: a C5 convertase.
    \item The C3b\({}_2\)Bb cleaves C5 into C5a and C5b. The C5b binds to the protein, attracting C6, C7, and C8 to bind on top.
    \item The C8 that binds on top allows for up to 16 individual C9 proteins to insert into the pathogen membrane forming a \href{https://en.wikipedia.org/wiki/Complement_membrane_attack_complex}{\textbf{membrane attack complex (MAC)}}.
\end{itemize}


\subsection{Classical pathway}

The classical pathway begins C1 binding to a pathogen that has already been tagged with appropriate antibodies. The C1 protein of three other proteins called C1q, C1r, and C1s. It is the C1s that cleaves C4 into C4a and C4b and the C4b that results binds to the pathogen. After, the C1s now cleaves C2 into C2a and C2b and the C2b binds onto the C4b forming C4b2b. This C4b2b is a C3 convertase.

\subsection{Alternative pathway}

Compared with the classical pathway, the alternative pathway does not require antibodies. It begins with the spontaneous cleavage of C3 into C3a and C3b.

\subsection{Lectin pathway}

The lectin pathway begins when \textbf{mannose-binding lectin (MBL)} or ficolins bind to characteristic carbohydrates on the surface of a pathogen. This activates \textbf{MBL-associated serine proteases (MASPs)}, principally MASP-2, which cleave C4 and C2. C4b binds to the pathogen surface and combines with C2b to form C4b2b: a C3 convertase.


\section{Cytokines and Chemokines}

\section{Antibiotics}

\begin{table}[htbp]
\centering
\footnotesize
\setlength{\tabcolsep}{3pt}
\renewcommand{\arraystretch}{1.2}

\begin{tabularx}{\textwidth}{
    @{}
    >{\RaggedRight\arraybackslash}p{3.5cm}
    >{\RaggedRight\arraybackslash}X
    @{}
}
    \toprule
    \textbf{Antibiotic} &
    \textbf{Notes} \\
    \midrule

    \multicolumn{2}{@{}l}{
        \textbf{\ensuremath{\boldsymbol{\beta}}-Lactam antibiotics}
        (inhibit cell-wall synthesis)
    } \\
    \cmidrule(lr){1-2}

    \hspace*{0.3cm}Phenoxymethylpenicillin &
    Penicillin used for infections such as streptococcal tonsillitis. \\

    \hspace*{0.3cm}Amoxicillin &
    Commonly used penicillin, particularly for respiratory infections. \\

    \hspace*{0.3cm}Flucloxacillin &
    Penicillin used particularly against susceptible
    \textit{Staphylococcus aureus}. \\

    \hspace*{0.3cm}Co-amoxiclav &
    Amoxicillin combined with the \ensuremath{\beta}-lactamase inhibitor
    clavulanic acid. \\

    \hspace*{0.3cm}Cefalexin &
    A cephalosporin used for selected respiratory, skin and urinary
    infections. \\


    \addlinespace[4pt]
    \multicolumn{2}{@{}l}{
        \textbf{Macrolides} (inhibit protein synthesis)
    } \\
    \cmidrule(lr){1-2}

    \hspace*{0.3cm}Clarithromycin &
    Used for respiratory infections and sometimes in penicillin allergy. \\

    \hspace*{0.3cm}Erythromycin &
    Similar to clarithromycin; also used for selected respiratory and
    skin infections. \\


    \addlinespace[4pt]
    \multicolumn{2}{@{}l}{
        \textbf{Tetracyclines} (inhibit protein synthesis)
    } \\
    \cmidrule(lr){1-2}

    \hspace*{0.3cm}Doxycycline &
    Used for respiratory infections, acne, chlamydia and some atypical
    infections. \\


    \addlinespace[4pt]
    \multicolumn{2}{@{}l}{
        \textbf{Aminoglycosides} (inhibit protein synthesis)
    } \\
    \cmidrule(lr){1-2}

    \hspace*{0.3cm}Gentamicin &
    Used mainly for serious Gram-negative infections; may cause renal
    and auditory toxicity. \\


    \addlinespace[4pt]
    \multicolumn{2}{@{}l}{
        \textbf{Fluoroquinolones} (inhibit DNA replication)
    } \\
    \cmidrule(lr){1-2}

    \hspace*{0.3cm}Ciprofloxacin &
    Used for selected Gram-negative infections; use is restricted because
    of potentially serious adverse effects. \\


    \addlinespace[4pt]
    \multicolumn{2}{@{}l}{
        \textbf{Nitroimidazoles} (damage bacterial DNA)
    } \\
    \cmidrule(lr){1-2}

    \hspace*{0.3cm}Metronidazole &
    Particularly active against obligate anaerobic bacteria. \\


    \addlinespace[4pt]
    \multicolumn{2}{@{}l}{
        \textbf{Antifolates} (inhibit folate metabolism)
    } \\
    \cmidrule(lr){1-2}

    \hspace*{0.3cm}Trimethoprim &
    Used for selected urinary tract infections. \\


    \addlinespace[4pt]
    \multicolumn{2}{@{}l}{
        \textbf{Nitrofurans}
    } \\
    \cmidrule(lr){1-2}

    \hspace*{0.3cm}Nitrofurantoin &
    Used mainly for lower urinary tract infections. \\


    \addlinespace[4pt]
    \multicolumn{2}{@{}l}{
        \textbf{Glycopeptides} (inhibit cell-wall synthesis)
    } \\
    \cmidrule(lr){1-2}

    \hspace*{0.3cm}Vancomycin &
    Used for serious Gram-positive infections, including some MRSA
    infections. \\

    \bottomrule
\end{tabularx}

\caption{Principal classes of antibiotics and selected examples }
\label{tab:antibiotics}

\end{table}



\begin{table}[htbp]
\centering
\footnotesize
\setlength{\tabcolsep}{3pt}
\renewcommand{\arraystretch}{1.2}

\begin{tabularx}{\textwidth}{
    @{}
    >{\RaggedRight\arraybackslash}p{2.9cm}
    >{\Centering\arraybackslash}p{1.8cm}
    >{\RaggedRight\arraybackslash}X
    @{}
}
    \toprule
    \textbf{Cell type} &
    \textbf{Phagocyte?} &
    \textbf{Notes} \\
    \midrule

    \multicolumn{3}{@{}l}{\textbf{Innate immune cells}} \\
    \cmidrule(lr){1-3}

    \hspace*{0.3cm}Neutrophils &
    Yes &
    Dead neutrophils are a major
    component of pus. \\

    \hspace*{0.3cm}Monocytes &
    Yes &
    Can differentiate into
    macrophages or monocyte-derived dendritic cells. \\

    \hspace*{0.3cm}Macrophages &
    Yes &
    \\

    \hspace*{0.3cm}Dendritic cells &
    Yes &
     \\

    \hspace*{0.3cm}Eosinophils &
    No &
     \\

    \hspace*{0.3cm}Basophils &
    No &
     \\

    \hspace*{0.3cm}Mast cells &
    No & Granulocytes; releases histamine.
    \\

    \hspace*{0.3cm}Natural killer (NK) cells &
    No &
    Recognise stressed, infected, or malignant cells by a lack of MHC, which triggers killing by apoptosis. \\

    \addlinespace[4pt]
    \multicolumn{3}{@{}l}{\textbf{Adaptive immune cells}} \\
    \cmidrule(lr){1-3}

    \hspace*{0.3cm}B lymphocytes &
    No &
     \\

    \hspace*{0.3cm}CD4\textsuperscript{+} helper T cells &
    No &
     \\

    \hspace*{0.3cm}CD8\textsuperscript{+} cytotoxic T cells &
    No &
    \\

    \hspace*{0.3cm}Regulatory T cells &
    No &
     \\

    \bottomrule
\end{tabularx}

\caption{Major cells of the innate and adaptive immune systems. Phagocytes are leukocytes specific to the immune system \cite[5:20]{CrashCourseImmune2}.}
\label{tab:immune-cells}

\end{table}


\subsection{What is the physiological use of inflammation?}

\begin{itemize}
    \item Swelling: vasodilated capillaries bring blood to the site of infection.
    \item Increased temperature: increases the cell metabolic rate to be able to reproduce faster. \cite[7:44]{CrashCourseImmune1}
\end{itemize}



\section{Bacteria}

\subsection{List of bacteria}
\begin{table}[htbp]
\centering
\footnotesize
\setlength{\tabcolsep}{3pt}
\renewcommand{\arraystretch}{1.2}

\begin{tabularx}{\textwidth}{
    @{}
    >{\RaggedRight\arraybackslash}p{3.5cm}
    >{\Centering\arraybackslash}p{1.0cm}
    >{\RaggedRight\arraybackslash}p{2.2cm}
    >{\RaggedRight\arraybackslash}X
    @{}
}
    \toprule
    \textbf{Organism/group} &
    \textbf{Gram} &
    \textbf{Shape} &
    \textbf{Example problems} \\
    \midrule

    \multicolumn{4}{@{}l}{\textbf{\textit{Staphylococcus}}} \\
    \cmidrule(lr){1-4}

    \hspace*{0.3cm}\textit{S. aureus} &
    $ + $ &
    Cocci in clusters &
    Skin infections, pneumonia, sepsis. \\

    \hspace*{0.3cm}\textit{S. epidermis} &
    ? &
    &
    Maybe less important \\

    % \hspace*{0.3cm}\textit{S. epidermidis} (coagulase-negative) &
    % $ + $ &
    % Cocci in clusters &
    % Line and prosthetic-device infections. \\

    \addlinespace[3pt]
    \multicolumn{4}{@{}l}{\textbf{\textit{Clostridioides}}} \\
    \cmidrule(lr){1-4}

    \hspace*{0.3cm}\textit{C. difficile} (C. diff) &
    $ + $ &
    Spore-forming rods &
    Antibiotic-associated diarrhoea and colitis. \\

    \hspace*{0.3cm}\textit{C. tetani}  &
    $ ? $ &
     &
    Causes tetanus \\

    \hspace*{0.3cm}\textit{C. perfringens}  &
    $ ? $ &
     &
     \\


    \addlinespace[3pt]
    \multicolumn{4}{@{}l}{\textbf{\textit{Pseudomonas}}} \\
    \cmidrule(lr){1-4}

    \hspace*{0.3cm}\textit{P. aeruginosa} &
    $ - $ &
    Rods &
    Pneumonia, wound infections and UTIs. Found in the environment, particularly moist areas. \\

    \addlinespace[3pt]
    \multicolumn{4}{@{}l}{\textbf{\textit{Enterococcus}}} \\
    \cmidrule(lr){1-4}

    \hspace*{0.3cm}\textit{E. faecalis} &
    $ + $ &
    Cocci in pairs or chains &
    UTIs, wound infections and endocarditis. \\

    \hspace*{0.3cm}\textit{E. faecium} &
    $ + $ &
    Cocci in pairs or chains &
    Healthcare-associated infections and bacteraemia. \\

    \addlinespace[3pt]
    \multicolumn{4}{@{}l}{\textbf{\textit{Streptococcus}}} \\
    \cmidrule(lr){1-4}

   
    % \hspace*{0.3cm}\textit{S. agalactiae} (group B) &
    % $ + $ &
    % Cocci in chains &
    % Neonatal sepsis and meningitis. \\

    \hspace*{0.3cm}\textit{S. pneumoniae} &
    $ + $ &
    Lancet-shaped diplococci &
    Pneumonia, otitis media and meningitis. \\

    \hspace*{0.3cm}\textit{S. mutans}  &
    $ + $ &
    Cocci in chains &
    Normal flora of pharynx and mouth.\\


    \addlinespace[3pt]
    \multicolumn{4}{@{}l}{\textbf{\textit{Escherichia}}} \\
    \cmidrule(lr){1-4}

    \hspace*{0.3cm}\textit{E. coli} &
    $ - $ &
    Rods &
    UTIs, diarrhoea and sepsis. \\

     \addlinespace[3pt]
    \multicolumn{4}{@{}l}{\textbf{\textit{Bacteroides}}} \\
    \cmidrule(lr){1-4}

    \hspace*{0.3cm}\textit{B. fragilis} &
    $ - $ &
    Rods & \\

     \addlinespace[3pt]
    \multicolumn{4}{@{}l}{\textbf{\textit{Candida}}} \\
    \cmidrule(lr){1-4}

    \hspace*{0.3cm}\textit{C. albicans} &
    NA &
    Yeast & \\

    \bottomrule
\end{tabularx}

\caption{Selected clinically important bacteria arranged by genus.
Bold italic headings indicate genera, while indented entries indicate
species. A `$ + $' indicates Gram-positive bacteria and a `$ - $'
Gram-negative bacteria. }
\label{tab:common-bacteria}

\end{table}

\subsection{Oxygen requirements}

Obligate anaerobes. Microaerobes are unique in that too much oxygen are toxic, but they require some oxygen. They won't grow in atmospheric oxygen concentrations, but they also don't grow in zero oxygen environments.

The less specific ones are the facultative anaerobes (usually aerobic, but \textit{can} be anaerobic) and facultative aerobes (vice versa)

Anaerobes are particularly important, because after an initial reaction, abscesses will become a low oxygen environment.


%The 
%
%
%Brain derived neurotropic factor
%
%BPA (Bisphenol A)
%very small amount toxic

